% GENERATED FILE. Edit canonical lessons, then run npm run generate:books.
% canonical-source-hash: fnv1a64:9447bdb9f88a315e
% canonical-lessons: GE-C483-gips, GE-C483-kruecke, GE-C483-humpeln, GE-C483-rollstuhl, GE-C483-heilen

\chapter{Gips und Krücken}
\label{ch:ge-gips-und-krucken}
\hlchaptermodality{483}

\begin{chapteropening}
\textbf{By the end of this chapter:} \emph{I can talk about broken bones and healing.}

Complete the last lesson of chapter 483: I can talk about broken bones and healing.
\end{chapteropening}

\section[der Gips]{der Gips — plaster cast; plaster}

\label{lesson:GE-C483-gips}



\begin{quote}
Before the new one: say the German for to disinfect, then the German for scar.
\end{quote}

\subsection*{You'll want to know: der Gips}
\textbf{der Gips} — \textquotedblleft{}plaster cast; plaster\textquotedblright{}.

\begin{quote}
Mein Arm war sechs Wochen im Gips.
\end{quote}

\begin{quote}
\emph{My arm was in plaster for six weeks.}
\end{quote}

\emph{Einen Gips bekommen} — to get a cast. Plural, rare: \emph{die Gipse}.

\begin{cousinweb}
\emph{Gips} comes from Middle High German \emph{gips}, from Latin \emph{gypsum}, which was borrowed from Greek \emph{gýpsos} \textquotedblleft{}gypsum\textquotedblright{} (EtymDB).
\end{cousinweb}

\begin{grammarlens}[title={What you've built}]
One more word for telling, explaining and describing at B1.
\end{grammarlens}

\subsection*{Your turn}
\begin{itemize}
\raggedright
  \item \emph{Say it:} \emph{der Gips}
  \item \emph{Say it:} \emph{der Gips}, once more
  \item \emph{Say it:} a sentence of your own with \emph{der Gips}
  \item \emph{Recall:} say the German for to disinfect, then the German for scar, then say \emph{der Gips} again
\end{itemize}

\begin{tcolorbox}[breakable,colback=teal!4,colframe=teal!35!black,title={Before you move on}]
What does der Gips mean? (\textquotedblleft{}Plaster cast; plaster\textquotedblright{}.) Say it once more.
\end{tcolorbox}
\section[die Krücke]{die Krücke — crutch}

\label{lesson:GE-C483-kruecke}



\begin{quote}
Before the new one: say the German for scar, then the German for plaster cast.
\end{quote}

\subsection*{You'll want to know: die Krücke}
\textbf{die Krücke} — \textquotedblleft{}crutch\textquotedblright{}.

\begin{quote}
Nach der Operation ging er an Krücken.
\end{quote}

\begin{quote}
\emph{After the operation he walked on crutches.}
\end{quote}

Plural: \emph{die Krücken}. \emph{An Krücken gehen} — to walk on crutches.

\begin{cousinweb}
\emph{Krücke} is inherited from Middle High German \emph{krucke} (EtymDB), a Germanic word; English \emph{crutch} is a cousin.
\end{cousinweb}

\begin{grammarlens}[title={What you've built}]
One more word for telling, explaining and describing at B1.
\end{grammarlens}

\subsection*{Your turn}
\begin{itemize}
\raggedright
  \item \emph{Say it:} \emph{die Krücke}
  \item \emph{Say it:} \emph{die Krücke}, once more
  \item \emph{Say it:} a sentence of your own with \emph{die Krücke}
  \item \emph{Recall:} say the German for scar, then the German for plaster cast, then say \emph{die Krücke} again
\end{itemize}

\begin{tcolorbox}[breakable,colback=teal!4,colframe=teal!35!black,title={Before you move on}]
What does die Krücke mean? (\textquotedblleft{}Crutch\textquotedblright{}.) Say it once more.
\end{tcolorbox}
\section[humpeln]{humpeln — to limp, to hobble}

\label{lesson:GE-C483-humpeln}



\begin{quote}
Before the new one: say the German for plaster cast, then the German for crutch.
\end{quote}

\subsection*{You'll want to know: humpeln}
\textbf{humpeln} — \textquotedblleft{}to limp, to hobble\textquotedblright{}.

\begin{quote}
Nach dem Spiel humpelte der Spieler vom Platz.
\end{quote}

\begin{quote}
\emph{After the game, the player limped off the pitch.}
\end{quote}

Perfekt with \emph{sein} when you go somewhere: \emph{Er ist nach Hause gehumpelt.}

\begin{cousinweb}
\emph{humpeln} is usually traced to Low German, and EtymDB sets it beside Dutch \emph{hompelen} \textquotedblleft{}to hobble\textquotedblright{}.
\end{cousinweb}

\begin{grammarlens}[title={What you've built}]
One more word for telling, explaining and describing at B1.
\end{grammarlens}

\subsection*{Your turn}
\begin{itemize}
\raggedright
  \item \emph{Say it:} \emph{humpeln}
  \item \emph{Say it:} \emph{humpeln}, once more
  \item \emph{Say it:} a sentence of your own with \emph{humpeln}
  \item \emph{Recall:} say the German for plaster cast, then the German for crutch, then say \emph{humpeln} again
\end{itemize}

\begin{tcolorbox}[breakable,colback=teal!4,colframe=teal!35!black,title={Before you move on}]
What does humpeln mean? (\textquotedblleft{}To limp, to hobble\textquotedblright{}.) Say it once more.
\end{tcolorbox}
\section[der Rollstuhl]{der Rollstuhl — wheelchair}

\label{lesson:GE-C483-rollstuhl}



\begin{quote}
Before the new one: say the German for crutch, then the German for to limp.
\end{quote}

\subsection*{You'll want to know: der Rollstuhl}
\textbf{der Rollstuhl} — \textquotedblleft{}wheelchair\textquotedblright{}.

\begin{quote}
Seit dem Unfall sitzt sie im Rollstuhl.
\end{quote}

\begin{quote}
\emph{Since the accident she has been in a wheelchair.}
\end{quote}

Plural: \emph{die Rollstühle}.

\begin{cousinweb}
\emph{Rollstuhl} is built on \emph{rollen} \textquotedblleft{}to roll\textquotedblright{} and \emph{Stuhl} \textquotedblleft{}chair\textquotedblright{} (EtymDB), which is inherited from Old High German \emph{stuol}, from Proto-Germanic \emph{stōlaz}.
\end{cousinweb}

\begin{grammarlens}[title={What you've built}]
One more word for telling, explaining and describing at B1.
\end{grammarlens}

\subsection*{Your turn}
\begin{itemize}
\raggedright
  \item \emph{Say it:} \emph{der Rollstuhl}
  \item \emph{Say it:} \emph{der Rollstuhl}, once more
  \item \emph{Say it:} a sentence of your own with \emph{der Rollstuhl}
  \item \emph{Recall:} say the German for crutch, then the German for to limp, then say \emph{der Rollstuhl} again
\end{itemize}

\begin{tcolorbox}[breakable,colback=teal!4,colframe=teal!35!black,title={Before you move on}]
What does der Rollstuhl mean? (\textquotedblleft{}Wheelchair\textquotedblright{}.) Say it once more.
\end{tcolorbox}
\section[heilen]{heilen — to heal}

\label{lesson:GE-C483-heilen}



\begin{quote}
Before the new one: say the German for to limp, then the German for wheelchair.
\end{quote}

\subsection*{You'll want to know: heilen}
\textbf{heilen} — \textquotedblleft{}to heal\textquotedblright{}.

\begin{quote}
Die Wunde heilt langsam.
\end{quote}

\begin{quote}
\emph{The wound is healing slowly.}
\end{quote}

Perfekt with \emph{sein} when a wound heals: \emph{Die Wunde ist gut geheilt.} With \emph{haben} when a doctor heals someone.

\begin{cousinweb}
\emph{heilen} is built on \emph{heil} \textquotedblleft{}whole, unhurt\textquotedblright{}, which is inherited from Old High German \emph{heil}, from Proto-Germanic \emph{hailaz} (EtymDB): to make whole.
\end{cousinweb}

\begin{grammarlens}[title={What you've built}]
One more word for telling, explaining and describing at B1.
\end{grammarlens}

\subsection*{Your turn}
\begin{itemize}
\raggedright
  \item \emph{Say it:} \emph{heilen}
  \item \emph{Say it:} \emph{heilen}, once more
  \item \emph{Say it:} a sentence of your own with \emph{heilen}
  \item \emph{Recall:} say the German for to limp, then the German for wheelchair, then say \emph{heilen} again
\end{itemize}

\begin{tcolorbox}[breakable,colback=teal!4,colframe=teal!35!black,title={Before you move on}]
What does heilen mean? (\textquotedblleft{}To heal\textquotedblright{}.) Say it once more.
\end{tcolorbox}
